\documentclass[lineno]{biometrika}

\usepackage{amsmath}
\usepackage{graphicx}
\usepackage{newtxtext}
\usepackage[subscriptcorrection]{newtxmath}
\usepackage{xr}
\externaldocument[S-]{supplement}

\def\Bka{{\it Biometrika}}
\newcommand{\pr}{\mathrm{pr}}
\newcommand{\LC}{\mathcal{L}}
\newcommand{\var}{\mathrm{var}}

\begin{document}

\jname{Biometrika}
\jyear{2026}
\jvol{}
\jnum{}
\cyear{2026}
\accessdate{}
%\received{}
%\revised{}

\markboth{K. W. Park}{Latent coverage from noisy calibration}

\title{Latent coverage from noisy calibration}

\author{KUN WOO PARK}
\affil{Department of Political Science and International Relations, Kookmin University,\\ 77 Jeongneung-ro, Seongbuk-gu, Seoul 02707, Republic of Korea \email{pkw31386094@gmail.com}}

\maketitle

\begin{abstract}
Prediction intervals for a latent quantity, such as the mean of an unsampled area or a response measured with error, are often calibrated on noisy proxies of that quantity. We ask what latent coverage a threshold calibrated on noisy scores guarantees when the law and the variance of the noise are unknown. If the latent distribution is bi-log-concave, a condition that allows asymmetry and some multimodality, noisy coverage $0.9$ guarantees latent coverage $0.8991$ under Gaussian noise, $0.8981$ under symmetric unimodal noise and $0.8935$ under mean-zero log-concave noise; without a shape assumption, median-zero noise guarantees only $0.8$. A unimodal latent distribution gives no more than the shape-free guarantee, and mean-zero unimodal noise allows no guarantee. When the noise laws and variances differ between units and are unknown, raising the rank of the noisy conformal threshold by a few places gives latent coverage with prescribed probability over calibration samples. We characterize the smallest such rank, exactly for symmetric unimodal noise. Without the raise, this probability can tend to zero as the calibration sample grows. Under Gaussian noise of known variance, the location of the latent mean determines whether the required widening of the radius is of the order of the noise variance or of its standard deviation.
\end{abstract}

\begin{keywords}
Conformal prediction; Log-concavity; Measurement error; Small area estimation; Tolerance interval.
\end{keywords}

\section{Introduction}\label{sec:intro}

Suppose that a residual $W$, the latent target minus a centre fitted on independent training data, is observed only through $V = W + e$, with noise $e$ independent of $W$. A survey statistician who calibrates a prediction interval for the true mean of an unsampled area on the direct estimates of sampled areas is in this position \citep{fay1979}, as is a conformal analyst whose calibration responses carry measurement error \citep{einbinder2024, cohen2025}. The threshold $t$ at which the noisy scores $|V|$ have coverage $p$ is observable. We ask what latent coverage $\pr(|W| \le t)$ it guarantees, uniformly over a class of latent laws and over noise laws of every scale.

Two answers are known. If $W$ and $e$ are both symmetric and unimodal, the theorem of \citet{anderson1955} gives $\pr(|W + e| \le t) \le \pr(|W| \le t)$, so the noisy threshold is conservative \citep{einbinder2024}. Without a shape assumption, median-zero noise turns a noisy level $1 - \alpha$ into a latent level $1 - 2\alpha$ \citep{guille2024}, and this cannot be improved.

Related methods target different quantities. LatentCP \citep{zheng2026} gives a marginal latent guarantee through a known forward model, and \citet{cohen2025} estimate the noise-free threshold by deconvolution, without a latent guarantee. Conformal intervals for small areas \citep{bersson2024} cover observable responses, and conformal confidence intervals \citep{zhang2026} target the expectations of sampled units. Empirical Bayes and parametric small-area intervals \citep{armstrong2022, ignatiadis2022, hallmaiti2006, yoshimori2014} concern sampled units or rely on a model for the latent effects.

We assume instead that the latent distribution is bi-log-concave \citep{dumbgen2017}. This class contains the log-concave laws and some multimodal ones, and does not require symmetry. Under it the loss in latent coverage is small, and we bound it for each class of noise laws (\S~\ref{sec:general}, Table~\ref{tab:map}). Unimodality of the latent law does not suffice, and neither does mean-zero unimodality of the noise. When the noise laws and variances differ between units and are unknown, the noisy conformal threshold at a slightly raised rank covers the latent target, and we characterize the smallest such rank (\S~\ref{sec:procedure}). For Gaussian noise of known variance, \S~\ref{sec:mean} studies the radius instead: the location of the latent mean decides whether the required widening is of the order of the noise variance, as for a fixed smooth law \citep{jochmans2024}, or of the noise standard deviation.

\section{Coverage transfer}\label{sec:general}

\subsection{One end of the interval}

Let $\LC$ be the class of log-concave laws on the real line, including point masses, and $\mathcal B \supset \LC$ the class of laws that are point masses or have a continuous distribution function $F$ with $\log F$ and $\log(1 - F)$ concave. Both classes are closed under translation and scaling, and convolution with a log-concave law preserves them \citep{prekopa1973}. Let $Q_q$ denote the lower $q$-quantile. Let $E_u$ be exponential with rate $u > 0$; we call the law of $b - E_u$, whose distribution function is $\min\{1, \mathrm{e}^{u(y - b)}\}$, an exponential tail. Let $\mathcal E$ be a class of noise laws closed under scaling and negation, and let
\[
\Psi_{\mathcal E}(q) = \sup_{\epsilon \in \mathcal E} E\{\min(1, q \mathrm{e}^{-\epsilon})\}.
\]

\begin{proposition}\label{prop:general}
Let $\epsilon$ have a law in $\mathcal E$ and be independent of $Y$, with $\log F_Y$ concave. If $\pr(Y + \epsilon \le 0) \ge p > \Psi_{\mathcal E}(q)$, then $F_Y(0) \ge q$. If $p < \Psi_{\mathcal E}(q)$, some exponential tail $Y$ with $F_Y(0) < q$ and some $\epsilon$ with law in $\mathcal E$ have $\pr(Y + \epsilon \le 0) > p$.
\end{proposition}

\begin{proof}
If $v = F_Y^{-1}(q) > 0$, the exponential tail $v + a/u - E_u$, $a = \log(1/q)$, $u$ a supergradient of $\log F_Y$ at $v$, lies stochastically below $Y$, and adding $\epsilon$ preserves this, so $p \le \pr(a/u - E_u + \epsilon \le 0) = E\{\min(1, q \mathrm{e}^{-u\epsilon})\} \le \Psi_{\mathcal E}(q)$, as $u\epsilon$ has a law in $\mathcal E$. If $F_Y$ has an atom at $v$ with $Y \ge v$, as for a point mass, then $p \le \pr(\epsilon < 0) \le \lim_{u\to\infty} E\{\min(1, q \mathrm{e}^{-u\epsilon})\} \le \Psi_{\mathcal E}(q)$. Conversely, $a - E_1$ with noise $\epsilon$ attains $E\{\min(1, q \mathrm{e}^{-\epsilon})\}$; shift it slightly to the right.
\end{proof}

Exponential tails are extremal because concavity of $\log F_Y$ limits how fast mass can accumulate below the $q$-quantile. Applied to $Y = (W - t)/s$ for any scale $s > 0$, Proposition~\ref{prop:general} controls the right end of $[-t, t]$; applied to $-W$, it controls the left end.

\subsection{Both ends}

A noisy failure probability $1 - p$ may be split in any way between the two ends, and each end transfers separately. Let $\lambda_{\mathcal E}(p) = \sup\{q' : \Psi_{\mathcal E}(q') < p\}$, the inverse of $\Psi_{\mathcal E}$ when $\Psi_{\mathcal E}$ is continuous and increasing, and $\phi_{\mathcal E}(a) = 1 - \lambda_{\mathcal E}(1 - a)$. By Proposition~\ref{prop:general}, an end with noisy failure probability $a$ has latent failure probability at most $\phi_{\mathcal E}(a)$. Define the sufficient noisy level
\begin{equation}\label{eq:pE}
p_{\mathcal E}(q) = \inf\{p : \phi_{\mathcal E}(a_L) + \phi_{\mathcal E}(a_R) \le 1 - q \text{ whenever } a_L + a_R \le 1 - p\}.
\end{equation}
Noisy coverage above $p_{\mathcal E}(q)$ at threshold $t$ then gives latent coverage $q$ for every $W \in \mathcal B$ and every noise with law in $\mathcal E$. Putting all the failure at one end, with small noise so that the other end is untouched, shows that no noisy level below $\Psi_{\mathcal E}(q)$ suffices (Supplementary \S~\ref{S-sec:s-oneend}). So $\Psi_{\mathcal E}(q) \le p_{\mathcal E}(q)$, with equality when $\phi_{\mathcal E}$ is convex, because the worst split then puts the whole failure at one end.

\begin{theorem}\label{thm:classes}
Let $q > \mathrm{e}^{-1}$.
(i) For symmetric unimodal noise, $\Psi_{\mathcal E}(q) = (1 + q \mathrm{e}^{-r})/2$, where $r > 0$ solves $(1 + r)\mathrm{e}^{-r} = (1 + \log q)/q$, attained by uniform noise; $\Psi_{\mathcal E}$ is strictly convex, so $\phi_{\mathcal E}$ is convex and $p_{\mathcal E}(q) = \Psi_{\mathcal E}(q)$.
(ii) For mean-zero log-concave noise, $\Psi_{\mathcal E}(q)$ is the supremum over centred log-affine laws on segments, and at $q = 0.9$ it lies in $[0.90517, 0.9052]$, the lower end given by the centred exponential law.
(iii) For mean-zero unimodal noise, $\Psi_{\mathcal E}(q) = 1$.
\end{theorem}

Part (i) holds because symmetric unimodal laws are scale mixtures of uniform laws; the proof of convexity is analytic (Supplementary Proposition~\ref{S-prop:su-exact}). For part (ii), the localization theorem \citep{fradelizi2004} reduces the supremum to log-affine laws on segments; a branch and bound in ball arithmetic bounds it, and a certified split of the failure probability between the ends gives $p_{\mathcal E}(0.9) \le 0.906$. For part (iii), a noise spread uniformly over $[-a, 0]$ with $a$ large, so that $q\mathrm{e}^{-\epsilon} \ge 1$ except on $[\log q, 0]$, and balanced to mean zero by a small mass far above zero, makes $E\{\min(1, q\mathrm{e}^{-\epsilon})\}$ arbitrarily close to $1$. For Gaussian noise, $p_{\mathcal E}(0.9) \in (0.90079, 0.901]$ (\S~\ref{sec:mean}).

\subsection{Unimodality does not suffice}

\begin{proposition}\label{prop:latentshape}
Let $\beta_{\mathcal E} = \inf_{\epsilon \in \mathcal E} \pr(\epsilon \ge 0) > 0$ and $1 - \beta_{\mathcal E} < p < 1$. Then $\pr(|W| \le 1) \ge 1 - (1 - p)/\beta_{\mathcal E}$ for every $W$ and every noise with law in $\mathcal E$ and $\pr(|W + e| \le 1) \ge p$. If the infimum is approached by laws continuous at $0$, then for every $\gamma > 0$ there are a unimodal $W$ with $E(W) = 0$ and such a noise with $\pr(|W| \le 1) \le 1 - (1 - p)/\beta_{\mathcal E} + \gamma$. If $\beta_{\mathcal E} = 0$, approached in the same way, then for every $p < 1$ and $\gamma > 0$ there are such $W$ and noise with $\pr(|W| \le 1) \le \gamma$.
\end{proposition}

For unimodal latent laws the guarantee is therefore the shape-free one. It is $2p - 1$ for median-zero noise, such as Gaussian or symmetric unimodal noise. It is $1 - \mathrm{e}(1 - p)$ for mean-zero log-concave noise, for which $\beta_{\mathcal E} = \mathrm{e}^{-1}$, attained by the centred exponential law. There is none for mean-zero unimodal noise, for which $\beta_{\mathcal E} = 0$.

The extremal laws place a thin layer of mass just outside the threshold, and the noise carries a fraction $1 - \beta_{\mathcal E}$ of it inside. Bi-log-concavity excludes such a layer: concavity of $\log F$ makes $f/F$ nonincreasing, so the density cannot jump up at the threshold. Proposition~\ref{prop:latentshape} shows that unimodality cannot replace bi-log-concavity; it does not show that bi-log-concavity is necessary.

\begin{table}
\tbl{Latent coverage guaranteed by noisy coverage $0.9$, the noisy level that guarantees latent coverage $0.9$, and the smallest valid rank (Theorem~\ref{prop:unknown}) at $\delta = 0.05$}
{\begin{tabular}{llcccc}
latent law & noise & latent coverage & noisy level & $k$, $K = 10^3$ & $k$, $K = 10^4$ \\
symmetric unimodal & symmetric unimodal & $\ge 0.9$ & $0.9$ & 916 & 9050 \\
$\mathcal B$ & Gaussian & $0.8991$ & $(0.90079, 0.901]$ & 917 & 9058--9060 \\
$\mathcal B$ & symmetric unimodal & $0.8981$ & $0.9017873\ldots$ & 918 & 9068 \\
$\mathcal B$ & mean-zero log-concave & $0.8935$ & $(0.90517, 0.906]$ & 921--922 & 9101--9109 \\
$\mathcal B$ & mean-zero unimodal & none & $1$ & -- & -- \\
unimodal, mean zero & median zero & $0.8$ & $0.95$ & 962 & 9537 \\
unimodal, mean zero & mean-zero log-concave & $1 - \mathrm{e}/10$ & $1 - 0.1/\mathrm{e}$ & 974 & 9664 \\
\end{tabular}}
\label{tab:map}
\begin{tabnote}
Coverages and levels are certified; a range of ranks reflects a level known within certified bounds, with the lower rank necessary and the upper sufficient. The first row is the usual high-probability rank, valid by Anderson's theorem; the exact level for symmetric unimodal noise is $\Psi_{\mathcal E}(0.9)$ of Theorem~\ref{thm:classes}(i), and the latent coverages are rounded down.
\end{tabnote}
\end{table}

\section{Calibration with unknown noise}\label{sec:procedure}

Let $W_1, \ldots, W_K, W_{\rm new}$ be independent with law $G$, and let $e_i = \sigma_i\epsilon_i$, independent of them and of each other. The laws of the $\epsilon_i$ lie in $\mathcal E$ and may differ between units, and the scales $\sigma_i > 0$ are unknown. The scores $V_i = W_i + e_i$ are independent but not identically distributed. Let $\mathcal D$ denote the calibration data and $T = |V|_{(k)}$. We call $\pr_{\mathcal D}\{\pr(|W_{\rm new}| \le T \mid \mathcal D) \ge q\}$ the latent reliability of rank $k$: the probability, over calibration samples, that the interval attains latent coverage $q$.

\begin{theorem}\label{prop:unknown}
If $G \in \mathcal B$, then
\[
\pr_{\mathcal D}\{\pr(|W_{\rm new}| \le T \mid \mathcal D) \ge q\} \ge \pr\{{\rm Bin}(K, p_{\mathcal E}(q)) \le k - 1\}.
\]
Conversely, if $\pr\{{\rm Bin}(K, \Psi_{\mathcal E}(q)) \le k - 1\} < 1 - \delta$, there are a log-concave $G$ and a noise law in $\mathcal E$, the same in every unit, for which the left side is below $1 - \delta$. Hence the smallest rank valid for every $G \in \mathcal B$ and all such noise lies between the smallest $k$ with $\pr\{{\rm Bin}(K, \Psi_{\mathcal E}(q)) \le k - 1\} \ge 1 - \delta$ and the smallest $k$ with $\pr\{{\rm Bin}(K, p_{\mathcal E}(q)) \le k - 1\} \ge 1 - \delta$, and equals both when $p_{\mathcal E}(q) = \Psi_{\mathcal E}(q)$.
\end{theorem}

\begin{proof}
Let $r = Q_q(|W|)$ for $W \sim G$. Latent coverage at $T$ is at least $q$ if and only if $T \ge r$, that is, if and only if $N = \#\{i : |V_i| < r\} \le k - 1$. For $s < r$, $\pr(|W| \le s) < q$, so by the transfer of \S~\ref{sec:general} at threshold $s$, $\pr(|V_i| \le s) \le p$ for every $p > p_{\mathcal E}(q)$, whatever the law and scale of $e_i$; letting $s \uparrow r$, $\pr(|V_i| < r) \le p_{\mathcal E}(q)$. $N$ is a sum of independent indicators with these success probabilities, so it is stochastically smaller than ${\rm Bin}\{K, p_{\mathcal E}(q)\}$. Conversely, for $p < \Psi_{\mathcal E}(q)$ the one-ended laws of \S~\ref{sec:general} give a log-concave $G$ and a noise with $\pr(|V_i| < r) \ge \pr(|V_i| \le 1) > p$ and $r > 1$, so the left side is at most $\pr\{{\rm Bin}(K, p) \le k - 1\}$, which tends to $\pr\{{\rm Bin}(K, \Psi_{\mathcal E}(q)) \le k - 1\}$ as $p \uparrow \Psi_{\mathcal E}(q)$.
\end{proof}

The bound is applied to each unit separately, so the noise laws and variances need not be known or shared. In practice: choose the noise class, read its noisy level from Table~\ref{tab:map}, and take as $k$ the smallest rank with $\pr\{{\rm Bin}(K, \text{level}) \le k - 1\} \ge 1 - \delta$. At $q = 0.9$ and $\delta = 0.05$ a valid rank exists for $K \ge 29$ under Gaussian or symmetric unimodal noise and for $K \ge 31$ under mean-zero log-concave noise; without a latent shape assumption and with median-zero noise it needs $K \ge 59$. For $K = 110$ all three bi-log-concave rows of Table~\ref{tab:map} give $k = 105$, the usual high-probability rank, against $109$ without a shape assumption, and the marginal rank $\lceil 0.9(K+1)\rceil = 100$ is not valid. The raise grows with $K$: the required level exceeds $q$ by a fixed amount, while consecutive ranks differ in level by about $1/K$. It cannot be dropped:

\begin{proposition}\label{prop:fail}
Let $q \in (0,1)$. There are a log-concave law $G$ and $D > 0$ such that, with $e_i \sim N(0, D)$ for all $i$, $\pr_{\mathcal D}\{\pr(|W_{\rm new}| \le |V|_{(k_K)} \mid \mathcal D) \ge q\} \to 0$ as $K \to \infty$ for every sequence of ranks with $k_K/K \to q$, such as the marginal rank and, for any fixed $\delta$, the usual high-probability rank.
\end{proposition}

\begin{proof}
By Theorem~\ref{thm:small} some log-concave $W$ has $H(1) \ge q$ and $r_q = Q_q(|W|) > 1$, where $H$ is the distribution function of $|W + D^{1/2}Z|$, $Z \sim N(0, 1)$. As in the proof of Theorem~\ref{prop:unknown}, the latent reliability of rank $k$ is $\pr\{{\rm Bin}(K, H(r_q-)) \le k - 1\}$, and $H(r_q-) > H(1) \ge q$ as $H$ is strictly increasing, so it tends to zero by the law of large numbers.
\end{proof}

A coverage deficit of $10^{-3}$ can thus remove the reliability. At the stress-test law of the Supplementary Material, with $D^{1/2}$ equal to $1\%$ of the noisy $q$-quantile, the latent coverage tends to $0.89918$. The usual rank has latent reliability $0.943$, $0.918$ and $0.151$ at $K = 10^3$, $10^4$ and $10^6$, against $0.954$, $0.958$ and $0.989$ for the rank of Theorem~\ref{prop:unknown} with Gaussian noise. For $29 \le K \le 300$ the usual rank's reliability is at most $0.005$ below the target $0.95$, so the raise matters more for conformal calibration on large sets of noisy regression responses than for small areas.

\section{Gaussian noise and the location of the mean}\label{sec:mean}

\subsection{The sharp radius}\label{sec:setup}

With Gaussian noise $e \sim N(0, D)$ of known variance, the noisy threshold can be widened instead of its rank raised, and we ask how much widening a guarantee requires. After scaling $t$ to $1$, write $x = D/t^2$, let $Z \sim N(0,1)$ be independent of $W$ and define
\begin{equation}\label{eq:R}
R_{p,q}(x) = \sup\{ Q_q(|W|) : W \in \LC,\ \pr(|W + x^{1/2}Z| \le 1) \ge p \},
\end{equation}
with $R_{p,q}(x) = -\infty$ if nothing is feasible. If calibration certifies noisy coverage $p$ at threshold $t$, then $t\,R_{p,q}(D/t^2)$ is the smallest radius that guarantees latent coverage $q$ for every log-concave latent law. The radius is sharp for this class and this single constraint; we do not claim optimality among procedures that use the whole calibration sample. The supremum may be restricted to point masses and laws with density proportional to $\mathrm{e}^{\beta w}$ on a segment, which makes $R_{p,q}$ computable (Supplementary Proposition~\ref{S-prop:reduction}). Let $R^{\mathcal B}_{p,q} \ge R_{p,q}$ be defined with $\mathcal B$ in place of $\LC$, and $R^{(b)}_{p,q}$ with the additional constraint $|E(W)| \le b$, so that the mean is at least $1 - b$ from $\pm 1$.

The behaviour of $R_{p,q}$ for small $x$ is governed by a one-sided constant,
\[
c_{p,q} = \sup\{F_Y^{-1}(q) : Y \in \LC,\ \pr(Y + Z \le 0) \ge p\}, \qquad c_q = c_{q,q},
\]
which is attained along exponential tails and is the same over $\mathcal B$ (Supplementary Lemma~\ref{S-lem:onesided}). Ball arithmetic gives $c_{0.8} = 0.04607$, $c_{0.9} = 0.01906$ and $c_{0.95} = 0.00841$.

\subsection{Three regimes}

Without a restriction on the mean the widening is of order $x^{1/2}$.

\begin{theorem}\label{thm:small}
For every $q \in (0,1)$, every $x > 0$ and every $W \in \mathcal B$ with $\pr(|W + x^{1/2}Z| \le 1) \ge q$, $Q_q(|W|) \le 1 + c_q x^{1/2}$. Moreover $R_{q,q}(x) = 1 + c_q x^{1/2} + o(x^{1/2})$ as $x \to 0$, and likewise for $R^{\mathcal B}_{q,q}$.
\end{theorem}

\begin{proof}[of the upper bound]
Let $\alpha_R = \pr(W + x^{1/2}Z > 1)$ and $\alpha_L = \pr(W + x^{1/2}Z < -1)$, so that $\alpha_L + \alpha_R \le 1 - q$. The law of $Y = (W - 1)/x^{1/2}$ has $\log F_Y$ concave and $\pr(Y + Z \le 0) = 1 - \alpha_R$. By Supplementary Lemma~\ref{S-lem:onesided}, which holds over $\mathcal B$ and gives that $c_q$ is nonincreasing in $q$, $F_Y^{-1}(1 - \alpha_R) \le c_{1-\alpha_R} \le c_q$, that is, $\pr(W > 1 + c_q x^{1/2}) \le \alpha_R$. The same argument applied to $-W$, which uses concavity of $\log(1 - F_W)$, gives $\pr(W < -1 - c_q x^{1/2}) \le \alpha_L$, so $\pr(|W| \le 1 + c_q x^{1/2}) \ge q$.
\end{proof}

The lower bound uses exponential tails whose right end lies just outside the threshold. Feasibility requires $x^{1/2} \le 1/\Phi^{-1}\{(1+q)/2\}$ (Supplementary \S~\ref{S-sec:s-feasible}), so at $q = 0.9$ the widening is below $1.16\%$ of the radius for every $W \in \mathcal B$, and below $0.22\%$ for log-concave laws by numerical upper bounds of $R_{q,q}$ (Supplementary Material).

A mean away from the boundary reduces the order to $x$.

\begin{theorem}\label{thm:mean}
Let $W \in \mathcal B$ with $|E(W)| \le B$, let $\sigma > 0$ and $t > B$, and suppose that $\pr(|W + \sigma Z| \le t) \ge q$. Then
\[
Q_q(|W|) \le B + \{(t - B)^2 + \sigma^2\}^{1/2} \le t + \frac{\sigma^2}{2(t - B)}.
\]
In particular $R^{(b)}_{p,q}(x) \le 1 + x/\{2(1 - b)\}$ for all $p \ge q$ and $x > 0$.
\end{theorem}

A mean far from an end of the interval caps how steeply the density can rise towards it. The proof makes this precise with the heat equation for quantiles \citep{steinbrecher2008}, which gives $d\{r(s) - B\}^2/ds \ge -1$ for $r(s) = Q_q(|W + s^{1/2}Z|)$. The order $x$ cannot be improved, even for mean-zero laws: centred log-affine laws on a segment attain it (Supplementary Proposition~\ref{S-prop:centred}), so centring does not replace symmetry.

The order $x^{1/2}$ returns only when the mean is within a bounded number of noise standard deviations of the boundary. Let
\[
L_q(\kappa) = \sup\{F_Y^{-1}(q) : Y \in \mathcal B,\ E(Y) \le -\kappa,\ \pr(Y + Z \le 0) \ge q\}.
\]

\begin{theorem}\label{thm:transition}
Let $1 - \mathrm{e}^{-1} < q < 1$, $x > 0$ and $W \in \mathcal B$ with $\pr(|W + x^{1/2}Z| \le 1) \ge q$. Then $|E(W)| < 1$, and if $\Delta = Q_q(|W|) - 1 > 0$,
\begin{equation}\label{eq:necessary}
2\{1 - |E(W)|\}\Delta + \Delta^2 \le x .
\end{equation}
Moreover, for every $\kappa \ge 0$, with $b_x = 1 - \kappa x^{1/2}$,
\[
\lim_{x \to 0} x^{-1/2}\{R^{(b_x)}_{q,q}(x) - 1\} = L_q(\kappa),
\]
and the same limit holds with $\LC$ replaced by $\mathcal B$ in the definition of $R^{(b)}$.
\end{theorem}

In the original units the mean lies $\kappa D^{1/2}$ from the threshold. The function $L_q$ equals $c_q$ up to a value $\kappa^*_q$, about $20$ at $q = 0.9$, and then decays like $\{1 - \log(1/q)\}/(2\kappa)$ (Supplementary Lemma~\ref{S-lem:meanonesided}).

\subsection{Slack and known variances}\label{sec:slack}

A noisy level $p$ slightly above $q$ removes the widening. Splitting the noisy and latent failure probabilities between the ends gives $Q_q(|W|) \le 1 + C_{p,q}x^{1/2}$ for every $x$ and every $W \in \mathcal B$ feasible at level $p \ge q$, where
\begin{equation}\label{eq:C}
C_{p,q} = \sup_{\alpha_L + \alpha_R = 1 - p}\ \inf_{\beta_L + \beta_R = 1 - q} \max(c_{1-\alpha_L, 1-\beta_L},\, c_{1-\alpha_R, 1-\beta_R}) \ge c_{p,q}.
\end{equation}

\begin{corollary}\label{cor:slack}
If $C_{p,q} \le 0$ then $R^{\mathcal B}_{p,q}(x) \le 1$ for every $x$, and if $c_{p,q} > 0$ then $R_{p,q}(x) > 1$ for all small $x$. At $q = 0.9$ the smallest slack $p - q$ with $C_{p,q} \le 0$, and with $R_{p,q}(x) \le 1$ for all $x$, lies in $(7.9\times10^{-4}, 10^{-3}]$. So for Gaussian noise $p_{\mathcal E}(0.9) \in (0.90079, 0.901]$, and noisy coverage $0.9$ gives latent coverage $0.8991$ for every $W \in \mathcal B$ and noise variance, while $0.8992$ fails for some log-concave laws.
\end{corollary}

\label{sec:hetero}Known variances allow a shorter interval. With heterogeneous known variances $D_i$, let $\bar g(w) = K^{-1}\sum_i \pr(|w + e_i| \le T)$ be the average noisy kernel at $T = |V|_{(k)}$. With probability $1 - \delta$ over calibration samples, $\int \bar g\, dG$ is at least the $\delta$-quantile $p_k$ of the ${\rm Beta}(k, K+1-k)$ distribution, when $k \ge Kp_k + 1$ (Supplementary Proposition~\ref{S-prop:beta}). The computed radius $T R^{\rm mix}$, with $R^{\rm mix}$ defined by \eqref{eq:R} under the constraint $\int \bar g\,dG \ge p_k - \omega$ and $\omega$ bounding a discretization error, then covers the latent target with probability $1 - \delta$ for log-concave $G$. This guarantee is for the exact $R^{\rm mix}$; our implementation computes a numerical upper bound on it, not a certificate. If only a lower bound $D_{\min}$ on the variances is known, the shrunk threshold $\max\{0, T + C_{p',q}D_{\min}^{1/2}\}$ has the same guarantee whenever $q \le p' \le p_k$ and $C_{p',q} < 0$; at $q = 0.9$, $K = 110$ and $\delta = 0.05$ this is $T - 0.114 D_{\min}^{1/2}$.

\section{Numerical illustration}\label{sec:numerical}

Table~\ref{tab:sameinfo} shows how much wider the noisy threshold must be to carry a latent guarantee when nothing is known about the noise. It compares the ranks of Theorem~\ref{prop:unknown} for several noise classes. No rule is told the noise law or the variances. The latent laws are normal, Laplace, centred gamma with shape $2$ and a truncated exponential law with density proportional to $\mathrm{e}^{4w}$ on $[0, 1]$, all standardized. The noise is Gaussian, uniform, Laplace, $t_3$ or centred exponential, with variances $D_i = 0.577\,L_i/E(L_i)$ and $L_i$ lognormal with log-scale $0.7$. Latent coverage given each data set is exact, and each of the $20$ settings has $2000$ data sets.

\begin{table}
\tbl{Rules with the same information: rank $k$, smallest reliability over the settings a rule covers, and extra width relative to the usual rank over the same settings, in percent; target reliability $0.95$}
{\begin{tabular}{lcccccc}
 & \multicolumn{3}{c}{$K = 110$} & \multicolumn{3}{c}{$K = 1000$} \\
rule & $k$ & reliability & width & $k$ & reliability & width \\
usual & 105 & -- & 0 & 916 & -- & 0 \\
Gaussian & 105 & 0.999 & 0 & 917 & 1.000 & 0.3--0.4 \\
symmetric unimodal & 105 & 0.998 & 0 & 918 & 1.000 & 0.6--1.0 \\
mean-zero log-concave & 105 & 0.999 & 0 & 922 & 1.000 & 2.0--2.9 \\
shape-free, median zero & 109 & 1.000 & 25--49 & 962 & 1.000 & 20--33 \\
shape-free, log-concave & 110 & 1.000 & 42--94 & 974 & 1.000 & 28--54 \\
\end{tabular}}
\label{tab:sameinfo}
\begin{tabnote}
Noisy levels $0.9$, $0.901$, $\Psi_{\mathcal E}(0.9)$, $0.906$, $0.95$ and $1 - 0.1/\mathrm{e}$. A rule covers a setting when the noise is in its class: Gaussian; symmetric unimodal (Gaussian, uniform, Laplace, $t_3$); mean-zero log-concave (Gaussian, uniform, Laplace, centred exponential); the first three need a bi-log-concave latent law. The usual rank covers no setting; its reliability over all $20$ settings was at least $0.998$.
\end{tabnote}
\end{table}

Every rule attained the target wherever it is guaranteed. The usual rank also did so here, because these laws are far from the extremal ones of Proposition~\ref{prop:fail}. At $K = 1000$ the shape-constrained ranks were at most $0.4\%$, $1.0\%$ and $2.9\%$ wider than the usual rank, and at $K = 110$ they coincide with it. The shape-free ranks were $20$--$54\%$ wider at $K = 1000$ and $25$--$94\%$ wider at $K = 110$. In these settings, assuming a bi-log-concave latent law reduced the extra width from $20$--$94\%$ to at most $2.9\%$.

Supplementary Table~\ref{S-tab:synth} compares rules that use the Gaussian noise variances. There the computed radius was $9.5$--$12\%$ shorter than the noisy threshold and reached the target for every law, and a Fay--Herriot tolerance interval adapted from \citet{chatterjee2008} was shorter still but missed the target under one log-concave law.

\section{Discussion}\label{sec:discussion}

The guarantees hold with probability $1 - \delta$ over calibration samples: on that event, the latent coverage of the new unit, given the calibration data, is at least $q$. They assume noise independent of the latent residual, and a latent law that is bi-log-concave, or log-concave for the computed radius, and shared by the calibration units and the new one. The target is an unsampled unit; for a sampled one the direct estimate is itself informative \citep{zhang2026}.

The guarantees allow heterogeneous noise within the specified classes, but they require a common latent distribution, and that assumption is the harder one to check. If the new residual is shifted by at most $\rho$ latent standard deviations and the latent law is log-concave, the latent coverage can fall by at most $\rho$, and by $q(1 - \mathrm{e}^{-\rho}) = q\rho + O(\rho^2)$ for some laws, so the linear order cannot be improved (Supplementary Proposition~\ref{S-prop:shift}). By comparison, at $q = 0.9$ Gaussian, symmetric unimodal and mean-zero log-concave noise cost at most $0.0065$ in coverage.

Bi-log-concavity cannot be checked on the latent scale. Convolution with log-concave noise preserves it, so with equal variances a confidence band for the signed noisy residuals that contains no bi-log-concave distribution function \citep{dumbgen2017} rules it out, though not conversely. Proposition~\ref{prop:latentshape} shows what the condition excludes: a thin layer of latent mass just beyond the threshold, which unimodality allows and which costs the full shape-free loss.

Open questions include an analytic proof that the centred exponential law is extremal among mean-zero log-concave noise laws, convexity of $\phi_{\mathcal E}$ for Gaussian and log-concave noise, which would make their two-sided levels exact as in Theorem~\ref{thm:classes}(i), and whether $R^{\mathcal B}_{p,q} = R_{p,q}$.

\section*{Declaration of the use of generative AI and AI-assisted technologies}

% To be written by the author.
[To be completed by the author.]

\section*{Supplementary material}\label{SM}

The Supplementary Material contains the omitted proofs, the certificates, further simulations and an application. Code reproducing every number is at \texttt{https://github.com/kunhailP/latent-coverage-noisy-thresholds}, tag \texttt{biometrika-submission}.

\bibliographystyle{biometrika}
\bibliography{refs}

\end{document}
